% ============================================================
%  Chapter 1 — Introduction
% ============================================================
\chapter{Introduction}
\label{cha:introduction}

% ------------------------------------------------------------
\section{Motivation}
\label{sec:motivation}

Consider a~patient in their sixties managing three chronic conditions: atrial fibrillation,
type~2 diabetes, and chronic kidney disease.
Three separate specialists (a~cardiologist, an endocrinologist, and a~nephrologist)
each prescribe their own regimen.
The result is a~daily combination of seven or more drugs, often without any single
physician having reviewed the complete medication list for cross-drug interactions.
Polypharmacy of this kind affects an estimated 40\% of patients over 65 in developed
countries~\cite{masnoon2017polypharmacy}, and drug--drug interactions (DDIs) contribute
to tens of thousands of preventable hospital admissions each year~\cite{lazarou1998ade}.

Modern large language models (LLMs) are impressively fluent and broad in their knowledge,
and they can, in principle, answer clinical questions about drug safety reasonably well.
But they are not safe for this task out of the box.
They hallucinate, their training knowledge has a~cutoff date, and they cannot
explain \emph{why} a~specific interaction exists or point to citable evidence for it.
Retrieval-Augmented Generation (RAG) was designed to address exactly these weaknesses:
by grounding every response in retrieved source documents, answers become traceable
and can be updated without retraining the model.

This thesis investigates how different RAG architectures perform on a~multi-hop
question-answering benchmark, proposes a~novel architecture tailored for high-stakes
advisory tasks, and demonstrates the approach through a~prototype drug interaction
advisory system.

% ------------------------------------------------------------
\section{Problem Statement}
\label{sec:problem}

Three research questions guide this work:

\begin{description}
  \item[RQ1] \label{rq:tradeoff}
    Which RAG architecture achieves the best trade-off between answer quality
    and latency, and does the answer depend on how hard the questions are?

  \item[RQ2] \label{rq:rare}
    Does a~routing-aware, set-wise evidence selection approach with a~grounding
    verifier (\rarerag{}) achieve higher faithfulness than single-path RAG architectures?

  \item[RQ3] \label{rq:platform}
    Can a~configurable microservice RAG platform serve as a~reliable technology
    demonstrator for drug interaction advisory, with results that are traceable to source
    documents?
\end{description}

RQ1 is posed with the difficulty clause deliberately.
An early version of this work asked simply which architecture is best, and the
evaluation in \cref{cha:evaluation} shows that the question has no
benchmark-independent answer: the architectures that lose on an encyclopedic
benchmark are the ones that win on hard clinical questions, using the same code.
The dependence on task difficulty appears to be the finding itself rather than a
complication of it.

Each research question is answered by a~specific experiment reported in a~specific
place, and \cref{tab:rq_map} states that mapping explicitly so that the empirical
basis of every answer can be located directly.

\begin{table}[htbp]
  \centering
  \caption[Mapping of research questions to experiments.]{Mapping of each research
    question to the experiment that verifies it, the section reporting it, and the
    answer obtained.  Every quality comparison is a~paired, per-question contrast
    against a~vanilla control, tested with a~Wilcoxon signed-rank test
    (\cref{subsec:stat_method}).}
  \label{tab:rq_map}
  \small
  \begin{tabularx}{\textwidth}{@{}l>{\raggedright\arraybackslash}X
                                  >{\raggedright\arraybackslash}X l@{}}
    \toprule
    \textbf{RQ} & \textbf{Experimental verification} & \textbf{Answer obtained}
      & \textbf{Section} \\
    \midrule
    RQ1
      & Nine pipeline variants on HotpotQA ($9 \times 1{,}000$) and on the DDI
        benchmark ($9 \times 100$, five difficulty levels); quality metrics and
        end-to-end latency, with a~per-difficulty breakdown
      & Benchmark-dependent: the agentic modes lose on HotpotQA and win on the
        hardest clinical questions.  Complexity buys grounding, not accuracy
      & \cref{sec:results_quality},
        \cref{subsec:ddi_difficulty},
        \cref{sec:results_latency} \\
    \addlinespace
    RQ2
      & \rarerag{} contrasted against Vanilla and the single-path modes on both
        benchmarks, with a~component-level diagnosis of the shortfall
      & Partially negative: the routing premise is supported, set-wise selection
        prunes evidence too aggressively, and \rarerag{} does not beat the vanilla
        baseline on either benchmark
      & \cref{subsec:rare_verdict},
        \cref{sec:hypothesis_eval} \\
    \addlinespace
    RQ3
      & End-to-end operation of the 14-service platform over the OpenFDA and
        DrugBank corpora, with citation traceability and production-mode metric
        collection; worked advisory queries
      & Yes as a~demonstrator: every claim is traceable to a~cited passage, with
        correctness near 0.5 and abstention never exercised by these benchmarks
      & \cref{sec:ddi_demo},
        \cref{cha:implementation} \\
    \bottomrule
  \end{tabularx}
\end{table}

% ------------------------------------------------------------
\section{Scope and Limitations}
\label{sec:scope}

The thesis covers a~survey of eight RAG architecture families (\cref{cha:background})
and the design of \rarerag{} as a~four-component pipeline (\cref{cha:rare_rag}).
A~14-service microservice platform with configurable RAG strategies is described
in \cref{cha:system_design,cha:implementation}, and evaluated empirically on
two benchmarks: HotpotQA ($9 \times 1{,}000$ questions for general multi-hop reasoning)
and a~purpose-built DDI benchmark ($9 \times 100$ questions at five clinical difficulty
levels), as reported in \cref{cha:evaluation}.

Several things are explicitly out of scope: clinical validation or patient-facing
deployment, fine-tuning or pre-training
of language models, and pharmacist-validated assessment of generated answers.

% ------------------------------------------------------------
\section{Contributions}
\label{sec:contributions}

The main contributions of this thesis are:

A~\textbf{comparative empirical study} of nine RAG pipeline variants on two benchmarks:
HotpotQA ($9 \times 1{,}000$ questions) and a~purpose-built DDI benchmark
($9 \times 100$ questions at five clinical difficulty levels), using lexical metrics
(Token~F1, EM, ROUGE-L), RAGAS-based indicators (faithfulness, answer correctness,
answer relevance) and latency.
Every comparison is a~paired, per-question contrast against a~vanilla control,
tested for significance, which turns out to matter, because several differences
that look decisive in the mode-level averages do not survive the test.

The central empirical finding: \textbf{the benefit of architectural complexity scales
with task difficulty, and it is paid in grounding rather than in accuracy}.
The four agentic architectures are significantly \emph{worse} than a~plain
retrieve-then-generate pass on HotpotQA, yet significantly better on the hardest
clinical questions, where a~single pass hallucinates on roughly one claim in two.
Crucially, this gain shows up in faithfulness and not in answer correctness:
across both benchmarks, no architecture makes the system measurably more likely to
be \emph{right}, only measurably more likely to be traceable to its sources.

The \textbf{\rarerag{} architecture} is a~RAG pipeline combining query complexity
routing, hybrid BM25+vector retrieval, set-wise evidence selection, and a~grounding
verifier with abstention.
It is reported here as a~partially negative result: the complexity router is
supported by the difficulty-scaling finding above, but set-wise selection prunes
evidence too aggressively on multi-drug questions, and \rarerag{} does not beat a
vanilla baseline on either benchmark (\cref{cha:evaluation}).

That negative result is treated here as a~contribution in its own right rather
than as a~shortfall to be minimised, and it is worth being explicit about why.
A~result of the form ``the proposed architecture wins'' is only as informative as
the baseline it was measured against.
What is reported instead is a~specific design error: optimising selected
evidence for diversity when the binding constraint on polypharmacy questions is
coverage, which is actionable for anyone building a~similar pipeline.
That precision is possible only because all nine modes run on one controlled
platform sharing a~single retrieval backend, embedding model and generator, so the
shortfall can be traced to the stage responsible instead of disappearing into an
end-to-end number.
Given how rarely unsuccessful configurations are reported at all, such localised
negative findings are correspondingly scarce.

The \textbf{open-source \medrag{} platform} is a~microservice system comprising
14 services with configurable chunking strategies, embedding providers, and RAG modes,
backed by a~full CI/CD pipeline and Docker deployment.

% ------------------------------------------------------------
\section{Thesis Organisation}
\label{sec:organisation}

\cref{cha:background} surveys the evolution of RAG from the original
retrieve-then-generate paradigm through agentic multi-agent systems and positions
this work within the existing literature.
\cref{cha:rare_rag} describes the \rarerag{} architecture in detail, covering
each of its four components and the hypotheses they are designed to test.
The \medrag{} platform's full microservice architecture is documented in
\cref{cha:system_design} using C4 diagrams and sequence models, and
\cref{cha:implementation} covers the most important engineering choices:
embedding abstraction, hybrid search tuning, SSE streaming, and evaluation integration.
Experimental results are presented and discussed in \cref{cha:evaluation}.
\cref{cha:conclusions} summarises what was found, revisits the research questions,
and suggests directions for future work.
